\documentclass{book}

\usepackage{fontspec}
\setmainfont{Ubuntu}

\usepackage[left=1in, right=1in, top=1in, bottom=1in]{geometry}

\usepackage{fancyhdr}
\pagestyle{fancy}
\lhead{
    \includegraphics[scale=0.05]{github.png}
    \href{https://github.com/QubitPi/general-relativity}{Study Notes}
}
\chead{General Relativity}
\rhead{
    \href{https://github.com/QubitPi}{QubitPi}
    \includegraphics[scale=0.05]{logo-8th-version.png}
}
\lfoot{}
\cfoot{}
\rfoot{\thepage}
\renewcommand{\headrulewidth}{0.4pt}
\renewcommand{\footrulewidth}{0.4pt}

\usepackage{hyperref}
\hypersetup{
    colorlinks=true,
    linkcolor=blue,
    anchorcolor=blue,
    urlcolor=blue
}

\usepackage{graphicx}
\usepackage{float}
\graphicspath{ {./img/} }

\usepackage{tikz}
\usetikzlibrary{calc, decorations.pathmorphing, quotes, angles, matrix, arrows.meta, hobby, decorations, decorations.markings}
\usepackage[most]{tcolorbox}
\tcbuselibrary{skins}
\tcbuselibrary{raster}

\usepackage{varioref}
\usepackage{cleveref}

\definecolor{myblue}{RGB}{0,163,243}
\NewTcbTheorem[auto counter, crefname={definition}{definitions}]
{Definition}{Definition}{
    breakable,
    parbox=false,
    fonttitle=\bfseries\upshape,
    fontupper=\slshape,
    arc=0mm,
    colback=myblue!20,
    colframe=myblue
}{def}

\newtcbtheorem[]{Theorem}{Theorem}
{
    breakable,
    parbox=false,
    theorem style=plain apart,
    label type=theorem,
    enhanced,
    frame hidden,
    boxrule=2mm,titlerule=0mm,toptitle=1mm,bottomtitle=1mm,
    fonttitle=\bfseries\large,fontupper=\normalsize,
    coltitle=green!35!black,colbacktitle=green!15!white,
    colback=green!50!yellow!15!white,
    borderline={1pt}{0pt}{green!25!blue},
}{theo}

\NewTcbTheorem[auto counter, crefname={corollary}{corollaries}]%
{Corollary}{Corollary}{
    breakable,
    parbox=false,
    colframe=green!50!black,colback=green!10!white,
    colbacktitle=green!20!white,coltitle=blue!75!black
}{cor}

\newtcbox{\roundinlinebox}[1][red]{
    on line,
    arc=7pt,
    colback=#1!10!white,
    colframe=#1!50!white,
    before upper={\rule[-3pt]{0pt}{10pt}},
    boxrule=1pt,
    boxsep=0pt,left=6pt,right=6pt,top=2pt,bottom=2pt
}

% https://tex.stackexchange.com/a/307436/277953
\newtcolorbox{marker}[1][]{enhanced,
before skip=2mm,after skip=3mm,
boxrule=0.4pt,left=5mm,right=2mm,top=1mm,bottom=1mm,
colback=yellow!50,
colframe=yellow!20!black,
sharp corners,rounded corners=southeast,arc is angular,arc=3mm,
underlay={%
    \path[fill=tcbcolback!80!black] ([yshift=3mm]interior.south east)--++(-0.4,-0.1)--++(0.1,-0.2);
    \path[draw=tcbcolframe,shorten <=-0.05mm,shorten >=-0.05mm] ([yshift=3mm]interior.south east)--++(-0.4,-0.1)--++(0.1,-0.2);
    \path[fill=yellow!50!black,draw=none] (interior.south west) rectangle node[white]{\Huge\bfseries !} ([xshift=4mm]interior.north west);
},
drop fuzzy shadow,#1}

\usepackage{varwidth}
\usepackage{amsfonts}
\usepackage{mathrsfs}
\usepackage{parskip}

\setlength{\parindent}{0pt}
\renewcommand{\labelitemi}{\includegraphics[width=0.05\textwidth]{嘟嘟可.png}}

% Text highlighting - https://tex.stackexchange.com/a/6029
\usepackage{soul}

\makeatletter

\newcommand{\defhighlighter}[3][]{%
    \tikzset{every highlighter/.style={color=#2, fill opacity=#3, #1}}%
}

\defhighlighter{yellow}{.5}

\newcommand{\highlight@DoHighlight}{
    \fill [ decoration = {random steps, amplitude=1pt, segment length=15pt}
        , outer sep = -15pt, inner sep = 0pt, decorate
        , every highlighter, this highlighter ]
    ($(begin highlight)+(0,8pt)$) rectangle ($(end highlight)+(0,-3pt)$) ;
}

\newcommand{\highlight@BeginHighlight}{
    \coordinate (begin highlight) at (0,0) ;
}

\newcommand{\highlight@EndHighlight}{
    \coordinate (end highlight) at (0,0) ;
}

\newdimen\highlight@previous
\newdimen\highlight@current

\DeclareRobustCommand*\highlight[1][]{%
    \tikzset{this highlighter/.style={#1}}%
    \SOUL@setup
%
    \def\SOUL@preamble{%
        \begin{tikzpicture}[overlay, remember picture]
            \highlight@BeginHighlight
            \highlight@EndHighlight
        \end{tikzpicture}%
    }%
%
    \def\SOUL@postamble{%
        \begin{tikzpicture}[overlay, remember picture]
            \highlight@EndHighlight
            \highlight@DoHighlight
        \end{tikzpicture}%
    }%
%
    \def\SOUL@everyhyphen{%
        \discretionary{%
            \SOUL@setkern\SOUL@hyphkern
            \SOUL@sethyphenchar
            \tikz[overlay, remember picture] \highlight@EndHighlight ;%
        }{%
        }{%
            \SOUL@setkern\SOUL@charkern
        }%
    }%
%
    \def\SOUL@everyexhyphen##1{%
        \SOUL@setkern\SOUL@hyphkern
        \hbox{##1}%
        \discretionary{%
            \tikz[overlay, remember picture] \highlight@EndHighlight ;%
        }{%
        }{%
            \SOUL@setkern\SOUL@charkern
        }%
    }%
%
    \def\SOUL@everysyllable{%
        \begin{tikzpicture}[overlay, remember picture]
            \path let \p0 = (begin highlight), \p1 = (0,0) in \pgfextra
            \global\highlight@previous=\y0
            \global\highlight@current =\y1
            \endpgfextra (0,0) ;
            \ifdim\highlight@current < \highlight@previous
            \highlight@DoHighlight
            \highlight@BeginHighlight
            \fi
        \end{tikzpicture}%
        \the\SOUL@syllable
        \tikz[overlay, remember picture] \highlight@EndHighlight ;%
    }%
    \SOUL@
}
\makeatother

\usepackage{amsthm}
\usepackage{cancel}
\usepackage{subcaption}
